\pdfbookmark[1]{Abstract}{Abstract}
\chapter*{Abstract}

A crystal carries an answer somebody else wrote down in advance. No other control in this work does.

Tile a published unit cell, carry every site as an exact integer, and score how well the arrangement
agrees with itself shifted along one axis. The offset where agreement peaks is the cell edge. Across
453 axes drawn from the Crystallography Open Database, every recovered period equals the published
edge as an integer, with no tolerance applied and no error left to average.

The same corpus read on a grid of 0.25 angstroms, with the tiling capped at 320 voxels along the
longest side, lands inside one voxel on every axis and on the published edge on none. Both bounds
come from the reader and neither comes from the deposit, and they cost a mean absolute error of
0.0124 angstroms together with a bias of $-0.0067$. The exact reading has a signed error of zero,
which places the bias entirely in the voxel floor. The cap also refuses large cells outright, a cost
that never appears as error because those entries go unread.

One correction here is not about the reader. A lattice of period $P$ agrees with itself as well at
$2P$ and at $3P$, and taking the tallest lag reports a harmonic: read that way, published edges come
back at almost exactly twice their value on ten of eighteen axes. Scoring a candidate together with
its multiples fixes it, and the rule holds with or without a grid, because it repairs a property of
periodic sets and not damage the reader does.

What this does not license is stated with it. The reading recovers the number the deposit carries
and stops there, only right-angled cells are used, and exact agreement does not make the
published edge true. The interest is not the cell edge, which is published to more decimal places
than anything here reaches. The interest is that a memoryless control can only show an instrument
does not invent structure, and a crystal lets it show that it finds structure that is there.
\cleardoublepage
